\documentclass{article}
\newcommand{\R}{\mathbb{R}}
\newcommand{\attn}[1]{\mathrm{Attn}(#1)}
\DeclareMathOperator{\softmax}{softmax}
\newtheorem{lem}{Lemma}

\begin{document}

\section{Introduction}
\label{sec:intro}
We study attention. % a comment with \cite{ignored}
The weights $A \in \R^{n \times n}$ are obtained via $\softmax$; see
Section~\ref{sec:method} and the result of \cite{vaswani2017,bahdanau2015}.

\section{Method}
\label{sec:method}
Given $d_k$, we scale by $\sqrt{d_k}$:
\begin{equation}
\label{eq:attn}
A = \softmax\left(\frac{QK^T}{\sqrt{d_k}}\right)
\end{equation}
By Lemma~\ref{lem:lipschitz}, using $d_k$ and $d_k$ again, with $W^Q$ and
$W^Q$ and $W^Q$ projections, the map is stable. Also $\alpha$, $\alpha$,
$\alpha$ appears often, as does $d_k$.

\begin{lem}[Stability]
\label{lem:lipschitz}
For any row-stochastic $A$, the mapping $A \mapsto AV$ is 1-Lipschitz.
\end{lem}

\begin{thebibliography}{9}
\bibitem{vaswani2017} Vaswani et al. Attention is all you need. NeurIPS 2017.
\bibitem{bahdanau2015} Bahdanau et al. Neural machine translation. ICLR 2015.
\end{thebibliography}

\end{document}
